% ---------------------------------------------------------------------------
% STUB. This section is NOT part of the draft: Section 3 (Preliminaries) is
% written later and will contain D1--D7 in full. The seven definitions below
% are copied verbatim from proofs/defs.tex so that the \cref labels used in
% Sections 4 and 5 resolve while those sections are being drafted.
% Replace this file by the real Section 3 and delete it.
% ---------------------------------------------------------------------------
\section{Preliminaries}\label{sec:prelim}

\begin{definition}[Field and domains]\label{def:domains}
$\F$ is a finite field; $k\ge 2$, $m\ge 1$ and $N$ are integers with $k^m\mid N$ and
$N\mid |\F|-1$. $L_0\subset\F^*$ is a coset of the subgroup of order $N$, and
$L_i:=\{x^k: x\in L_{i-1}\}$ for $1\le i\le m$, so $|L_i|=N/k^i$. The \emph{fibre} of
$y\in L_i$ is $\fib(y):=\{x\in L_{i-1}: x^k=y\}$; it equals $x_0\mu_k$ for any
$x_0\in\fib(y)$, where $\mu_k$ is the group of $k$-th roots of unity, and $|\fib(y)|=k$.
\end{definition}

\begin{definition}[Polynomials]\label{def:poly}
$R$ denotes a commutative $\F$-algebra (in practice $R=\F$ or $R=\F[T]$).
$R[x_i,\dots,x_m]_{<k}$ is the set of polynomials with $\deg_{x_l}<k$ for all $l$.
\end{definition}

\begin{definition}[Kronecker map]\label{def:kron}
For $0\le i\le m$ and $g\in R[x_{i+1},\dots,x_m]_{<k}$,
\[
U^{(i)}_g(Y):=g\big(Y,Y^{k},\dots,Y^{k^{m-i-1}}\big)\in R[Y],\qquad U_f:=U^{(0)}_f .
\]
For $i=m$, $g$ is a constant and $U^{(m)}_g=g$.
\end{definition}

\begin{definition}[Fold]\label{def:fold}
For $u:L_{i-1}\to R$ and $y\in L_i$, the \emph{components} of $u$ are
\[
u_j(y):=\frac1k\sum_{x\in\fib(y)}x^{-j}\,u(x)\qquad(0\le j<k),
\]
and for $c\in R$ the \emph{fold} is $\Fold_c(u):L_i\to R$,
$\Fold_c(u)(y):=\sum_{j<k}c^j\,u_j(y)$.
\end{definition}

\begin{definition}[Weighted and graded degree]\label{def:wdeg}
Let $w\in\N^m$. For $f\in\F[x_1,\dots,x_m]$,
$\deg_w f:=\max\{\langle w,\alpha\rangle: f_\alpha\ne0\}$.
For $g\in\F[T][x_{i+1},\dots,x_m]$, $\gdeg_w g:=\max\{a+\sum_{l>i}w_l\alpha_l\}$
over the monomials $T^a x^\alpha$ of $g$.
\end{definition}

\begin{definition}[Codes]\label{def:codes}
$\RS[\F,L,d]:=\{P|_L:\deg P<d\}$. The $t$-interleaved code $\RS[\F,L,d]^t$ is the set
of words $L\to\F^t$ whose $t$ coordinates lie in $\RS[\F,L,d]$; we identify
$\F[T]_{<t}$ with $\F^t$ through coefficients. The \emph{graded codes} are
\[
\Vc_i^{w,D}:=\big\{\,U^{(i)}_g\big|_{L_i}\ :\ g\in\F[T][x_{i+1},\dots,x_m]_{<k},\ \gdeg_w g\le D\,\big\},
\]
and the \emph{degree code} $C_{w,D}$ is the set of $T$-free elements of $\Vc_0^{w,D}$.
\end{definition}

\begin{definition}[Weighted agreement]\label{def:agree}
For words $u,v$ on $L$ and $\mu:L\to[0,1]$,
$\agree_\mu(u,v):=\frac1{|L|}\sum_{x:\,u(x)=v(x)}\mu(x)$ (as in BCIKS20, \S7).
\end{definition}
